\begin{center}
    {\large Міністерство освіти і науки України}\\[1em]
    {\large НАЦІОНАЛЬНИЙ УНІВЕРСИТЕТ ``КИЄВО-МОГИЛЯНСЬКА АКАДЕМІЯ''}\\[0em]
    {\large Кафедра мультимедійних систем\\факультету інформатики}\\[2em]

    {\large ЗАТВЕРДЖУЮ}\\
    {\large Зав. кафедри мультимедійних систем,}\\
    {\large доцент, кандидат наук,}\\
    {\large \underline{\hspace{3cm}} Жежерун О.П}\\
    {\mbox{\large (підпис) \hspace{4cm}}}\\
    ``\underline{\hspace{1cm}}'' \underline{\hspace{3cm}} 2025~р.\\[2em]
    
    {\large ІНДИВІДУАЛЬНЕ ЗАВДАННЯ}\\[0.5em]
    {\large на курсову роботу}\\[0.5em]
    студенту Гнутову Олександру 3 курсу\\факультету інформатики\\[0.5em]
    
    ТЕМА: Трасування променів у воксельному просторі\\[1em]
    
    \begin{flushleft}
        Зміст ТЧ до курсової роботи:
        \begin{itemize}[itemsep=0pt, topsep=0pt, label={}, leftmargin=18pt]
            \item Індивідуальне завдання
            \item Вступ
            \item \begin{enumerate}
                \item Аналіз предметної області та постановка задачі
                \item Методи та особливості реалізації трасування у воксельному просторі
                \item Реалізація алгоритму трасування променів у воксельному просторі
            \end{enumerate}
            \item Висновок
            \item Список використаних джерел
            \item Додатки
        \end{itemize}
        \vspace{1em}

        Дата видачі ``\underline{\hspace{1cm}}'' \underline{\hspace{3cm}} 2026~р.\\
        Керівник \underline{\hspace{5cm}} Завдання отримав: \underline{\hspace{3cm}}
    \end{flushleft}
\end{center}
\newpage